\item In the arrangement shown in Fig. 1.17 the mass of the rod \( M \) exceeds the mass \( m \) of the ball. The ball has an opening permitting it to slide along the thread with some friction. 
    \begin{center}
        \begin{tikzpicture}
            \node at (0, 0) {{f117.png}};
        \end{tikzpicture}
    \end{center}
    The mass of the pulley and the friction in its axle are negligible. At the initial moment the ball was located opposite the lower end of the rod. When set free, both bodies began moving with constant accelerations. Find the friction force between the ball and the thread at \( t \) seconds after the beginning of motion the ball got opposite the upper end of the rod. The rod length equals \( l \).